\documentclass[12pt]{article}

\usepackage[margin=1in]{geometry}
\usepackage{amsmath,amssymb,amsthm}
\usepackage[colorlinks=true, allcolors=blue]{hyperref}

\title{A Faster Way to Compute $c_t(n)$}
\author{Samuel Orellana Mateo\thanks{Duke University, Durham, NC 27708, USA. \href{mailto:samuelorellanamateo@gmail.com}{\texttt{samuelorellanamateo@gmail.com}}}}
\date{March 29, 2026}

\begin{document}

\maketitle

\begin{abstract}
    The function $c_t(n)$ is central to the enumeration of prime juggling patterns. This writeup provides a closed-form generating function for $c_t(n)$, enabling a faster algorithm to compute its values.
\end{abstract}

\vspace{1em}

The function $c_t(n)$ was originally introduced by Banaian et al.\ \cite{Banaian2016} to establish the exact count of normal 2-ball prime juggling patterns. It was subsequently utilized as a building block to enumerate prime patterns in multiplex and colored juggling variations \cite{butler2026}. By definition, $c_t(n)$ involves a sum over all partitions of $n$ into $t$ distinct parts $p_1 > p_2 > \dots > p_t \ge 1$:
\[
    c_t(n) = \sum_{\substack{p_1>\cdots>p_t \ge 1 \\ p_1+\cdots+p_t=n}}
    \frac{1}{t(t+1)}\prod_{i=1}^t\bigg(\frac{i+1}{i}\bigg)^{p_i}
\]

To obtain a generating function, we change variables from the partition parts $p_i$ to their adjacent differences. Let $\delta_t = p_t$, and $\delta_{i} = p_{i} - p_{i+1}$ for $i = 1, \dots, t-1$. It follows that $\delta_j \ge 1$ for all $1 \le j \le t$. We can write $p_i = \sum_{j=i}^t \delta_j$. The partition sum condition then becomes:
\[
    \sum_{i=1}^t p_i = \sum_{j=1}^t j \delta_j = n
\]

Substituting this into the product term gives a telescoping product:
\begin{align*}
    \prod_{i=1}^t \left(\frac{i+1}{i}\right)^{p_i} 
    &= \prod_{i=1}^t \left(\frac{i+1}{i}\right)^{\sum_{j=i}^t \delta_j} \\
    &= \prod_{j=1}^t \prod_{i=1}^j \left(\frac{i+1}{i}\right)^{\delta_j} \\
    &= \prod_{j=1}^t (j+1)^{\delta_j}
\end{align*}

Thus, $c_t(n)$ can be rewritten as a sum over these independent differences:
\[
    c_t(n) = \frac{1}{t(t+1)} \sum_{\substack{\delta_1, \dots, \delta_t \ge 1 \\ \sum_{j=1}^t j \delta_j = n}} \prod_{j=1}^t (j+1)^{\delta_j}
\]

Let $x$ be a formal variable. The constraint $\sum_{j=1}^t j \delta_j = n$ with the weight $(j+1)^{\delta_j}$ describe the coefficient of $x^n$ in the product of infinite geometric series. Let $G_t(x) = \sum_{n=1}^\infty c_t(n) x^n$ be the generating function for a fixed $t$. Then:
\[
    G_t(x) = \frac{1}{t(t+1)} \prod_{j=1}^t \sum_{\delta=1}^\infty \Big( (j+1) x^j \Big)^\delta
\]

Evaluating the $\sum_{\delta=1}^\infty r^\delta = \frac{r}{1-r}$, we obtain the closed-form:
\[
    G_t(x) = \frac{1}{t(t+1)} \prod_{j=1}^t \frac{(j+1)x^j}{1 - (j+1)x^j}
\]

To compute $c_t(n)$ algorithmically, one can just compute the coefficient $[x^n] G_t(x)$. 

Furthermore, we can extract a linear recurrence from this generating function to avoid the overhead of symbolic polynomial expansion. Let $H_t(x) = t(t+1) G_t(x)$, so that:
\[
    H_t(x) = \prod_{j=1}^t \frac{(j+1)x^j}{1 - (j+1)x^j}
\]
We can define $H_t(x)$ recursively:
\[
    H_t(x) = H_{t-1}(x) \frac{(t+1)x^t}{1 - (t+1)x^t}
\]
Multiplying both sides by the denominator gives:
\[
    H_t(x) \big(1 - (t+1)x^t\big) = H_{t-1}(x) (t+1)x^t
\]
By extracting the coefficient of $x^n$, which we denote as $h_t(n) = [x^n]H_t(x)$, this becomes:
\[
    h_t(n) = (t+1) \Big[ h_{t-1}(n-t) + h_t(n-t) \Big]
\]
with the base case $h_0(0) = 1$ and all other initial values set to $0$. The final count is then recovered by $c_t(n) = \frac{h_t(n)}{t(t+1)}$. This reduces the computation to a $O(n \sqrt{n})$ dynamic programming algorithm that relies only on integer addition and multiplication.

\vspace{2em}

\begin{thebibliography}{9}

\bibitem{Banaian2016}
Esther Banaian, Steve Butler, Christopher Cox, Jeffrey Davis, Jacob Landgraf, and Scarlitte Ponce.
\newblock Counting Prime Juggling Patterns.
\newblock \emph{Graphs and Combinatorics}, 32(5):1675--1688, Sep 2016.
\newblock \href{https://doi.org/10.1007/s00373-016-1711-1}{doi:10.1007/s00373-016-1711-1}.

\bibitem{butler2026}
Steve Butler, Vera Choi, Joel Jeffries, Nina McCambridge, Asia Morgenstern, and Samuel Orellana Mateo.
\newblock Enumerating Prime Patterns in Juggling Variations.
\newblock \emph{arXiv preprint arXiv:2603.17284}, 2026.
\newblock \url{https://arxiv.org/abs/2603.17284}.

\end{thebibliography}

\end{document}